\begin{textsample}{POS Dim 1 – ba – Score 47.00 – ba/ba\_em\_33.txt}  \label{ex:f1_pos_110}
SOCIOLOGIA A Sociologia é um componente curricular que compõe a área do conhecimento das \textbf{ciências} humanas, nasceu entre o \textbf{final} do \textbf{século} XIX e início do \textbf{século} XX e, por isso, sua história e trajetória são marcadas pelo desenvolvimento da sociedade \textbf{moderna}, urbana e industrial.
Deste modo, a transformação da vida social em objeto sociológico é um fenômeno da modernidade e a produção dos saberes decorrente deste fenômeno deve ser compreendida, portanto, como uma \textbf{demanda} da própria sociedade \textbf{moderna} secularizada.
É possível estabelecer tendências e interesses distintos para cada corrente de pensamento incluída na história da Sociologia.
Tais distinções correspondem a propósitos específicos de cada \textbf{autor} ou perspectiva \textbf{teórica}, sendo possível classificar as diversas interpretações de acordo com os objetivos ora ocultos, ora expressos das \textbf{teorias}.
Não se pretende propor um currículo que seja seletivo quanto à inclusão ou exclusão de \textbf{teorias}, interpretações e pesquisas sociológicas a partir de suas manifestações em espectros políticos.
Contudo, a neutralidade axiológica já foi definida pelo clássico Max Weber como sendo um exercício intelectual \textbf{ideal} típico necessário ao fazer científico, mas não como uma condição para o sociólogo, uma vez que as subjetividades dos agentes de uma pesquisa interagem entre si, fundindo observador e observado.
Assim, reconhecemos que a funcionalidade do saber sociológico está relacionada à \textbf{capacidade} de aplicação que lhe foi legada, quando os \textbf{teóricos} clássicos conceberam o que estudar, o que compreender e o que fazer a partir da \textbf{disciplina}.
Em outras palavras, a Sociologia \textbf{inspira} perícia na \textbf{análise} dos fatores da \textbf{coesão} social se lemos Durkheim, refaz o que tomamos por sociedade enquanto realidade objetiva quando lemos Weber e integra o que ficou à margem no processo histórico de alienação quando interpretamos Marx porque a relação que une os objetivos políticos ( não necessariamente partidários ) aos cognitivos de dada interpretação sociológica é o que comprova a própria \textbf{aplicabilidade} desta \textbf{disciplina}.
O recorte de um olhar denuncia a própria “ natureza ” de ser particular que é a realidade social.
É que a \textbf{razão} de ser do aprendizado na Sociologia reside na \textbf{capacidade} desta \textbf{ciência} de \textbf{instigar} a reflexão sobre as práticas sociais.
A consciência sobre a reificação do mundo torna -se, assim, o produto de uma consciência sociológica capaz de desnaturalizar os processos históricos e sociais.
Neste sentido, o produto do conhecimento sociológico alimenta a construção de uma cidadania agente, amparado e fundamentado numa legislação \textbf{comprometida} com o caráter humanista e progressista da formação na Educação Básica.
A história da Sociologia no currículo do Ensino Médio no Brasil inicia com a Lei n{\EmojiFont °}9.394/96, em seu Art. 36, § 1{\EmojiFont °}, III, que dispõe sobre a necessidade dos conhecimentos de Filosofia e Sociologia para a formação cidadã, e se consolida com Lei n{\EmojiFont °} 11.684/ 2008 que tornou obrigatório o ensino de Sociologia e Filosofia no Ensino Médio.
As \textbf{demandas} da \textbf{contemporaneidade}, em relação às premissas humanistas que norteiam valores éticos e morais para a construção de uma sociedade justa e inclusiva, colocam para a Sociologia a \textbf{tarefa} de discutir modelos \textbf{econômicos} desgastados, padrões de produção e consumo inviáveis para a conservação de recursos naturais e humanos, bem como outras questões relativas ao constante fluxo de transformações do mundo social.
Esta discussão tem, portanto, como principal objetivo garantir na formação de estudantes do Ensino Médio \textbf{compreensões} da força das macroestruturas sobre a dinâmica social ; perceber a historicidade e reversibilidade das próprias estruturas sociais ; avaliar o mundo que o cerca com um olhar \textbf{crítico}, assumindo sua subjetividade e \textbf{demandas} próprias.
O esforço para garantir esses objetivos de aprendizagem se traduziu no \textbf{organizador} curricular do componente, preservando referências aos clássicos da \textbf{disciplina}, \textbf{agregando} contribuições \textbf{teóricas}, \textbf{conceituais} e empíricas para o \textbf{debate} sociológico significativo e \textbf{relevante} para o tempo \textbf{atual}.
A tendência ora tomada pelo Novo Ensino Médio na Rede Estadual de Ensino da Bahia está alinhada à Base Nacional Comum Curricular, ainda que a transcenda propondo que a flexibilização do currículo não subtraia a aquisição do conhecimento historicamente acumulado.
A escolha da pedagogia histórico-crítica, cujo referencial é a pedagogia de Dermeval Saviani, se dá pela busca da construção das mudanças do currículo referencial da Bahia com um \textbf{horizonte} para a educação em seu sentido integral, considerando as nuances do processo de formação escolar e de construção do conhecimento, de modo que, o \textbf{teórico} ganhe sentido cognitivo uma vez, mostrando -se prático, vivencial e com significado para o estudante em sua jornada de formação e de inserção na sociedade.
As habilidades e \textbf{competências} organizadas pela BNCC enunciam a aquisição de conhecimentos cujo domínio perpassa o estudo de objetos sociológicos, \textbf{tanto} no âmbito do mundo do trabalho, quanto no universo sociocultural ou do campo das relações sociais de interação e das instituições e papéis sociais.
A metodologia do componente de Sociologia para o Ensino Médio deve se apropriar dos recursos disponíveis para referir às múltiplas linguagens que \textbf{figuram} atualmente no processo didático, de modo que permitam ao estudante participar do ensino-aprendizagem com protagonismo, diante das novas \textbf{tecnologias} educacionais.
Considerando a diversidade temática e \textbf{teórica} da \textbf{disciplina}, os \textbf{métodos} de \textbf{abordagem} dos conceitos e interpretações devem transparecer o caráter polifônico do discurso sociológico, \textbf{tanto} quanto se mostra diversa étnica-racial, sexual, cultural, geracional e regionalmente a extensão do corpo estudantil ao qual atende a Rede Estadual.
Por isso, o exercício da leitura \textbf{crítica} e \textbf{dialógica} fundamentam o tipo de \textbf{abordagem} aqui recomendada.
As \textbf{sequências} didáticas devem traduzir a distribuição do saber sociológico - em seus campos de pesquisa, marcos \textbf{teóricos} e metodologias - de forma contextualizada e correspondente às transformações históricas da \textbf{disciplina}, bem como da própria sociedade.
Contudo, não se deve tomar a divisão projetada como refratária aos possíveis \textbf{desdobramentos} do contexto vivido, pois sempre que necessário haverá espaço para integrar uma mesma discussão com base em mais de um referencial \textbf{teórico} e em todos os níveis do Ensino Médio, desde que o recorte proposto seja \textbf{pertinente} para a formação da cidadania e fortalecimento da democracia.
Assim, para cumprir objetivos como o de familiarizar o estudante com \textbf{problemas} sociológicos, propõe -se exercitar a \textbf{análise} da realidade social, oportunizando o \textbf{debate} interdisciplinar e estimulando a leitura e a produção textual, nas mais variadas linguagens.
A avaliação precisa estar atenta às metodologias ativas.
Deste modo, os critérios \textbf{avaliativos} observam a aquisição de habilidades como compreender as particularidades dos fenômenos humanos e de suas determinações, demonstrar autonomia para formular um olhar \textbf{crítico}, dominar os significados de conceitos políticos chave para a construção da cidadania e da democracia, distinguir e avaliar papéis referentes aos distintos \textbf{atores} sociais e políticos, conectar aspectos culturais e ideológicos às perspectivas de realidade social e de interesses políticos, compreender as lógicas que orientam as novas configurações sociais, atentando para as transformações do mundo do trabalho, da relação espaço-tempo, da família e de outras instituições sociais.
A escola \textbf{segue} sendo desafiada a assumir o papel de educar, a despeito do efeito contrastante que o cerceamento \textbf{econômico}, cultural e político produz sobre a Educação.
Bourdieau e Passeron ( 2018 ) demarcaram o campo da Sociologia da Educação quando observaram a herança cultural como fator de hierarquização social e o sistema de ensino modulado por essas distinções de status cultural, reproduzindo as desigualdades na escola e na universidade.
O diálogo que foi aberto em torno da instituição escola alinhou em um mesmo nível de importância a discussão, \textbf{tanto} pedagógica sobre as didáticas de ensino e os conteúdos curriculares, quanto sócio antropológica sobre as implicações culturais, políticas e \textbf{econômicas} da dinâmica escolar de integração social.
Esse diálogo entre o pedagógico e o sociológico, no papel da escola aponta para uma vontade de superar o caráter de reprodução das desigualdades sociais, para a vontade de construir uma escola democrática com seu acesso universal.
A escola precisa ser inclusiva, atender a especificidades plurais, integrar e acolher a diversidade com o espírito da alteridade.
É isso que o universalismo do projeto democrático de sociedade e de Educação \textbf{exige} ao ensino de Sociologia.

% matched lemmas: abordagem, agregar, análise, aplicabilidade, ator, atual, autor, avaliativo, capacidade, ciência, coesão, competência, compreensão, comprometido, conceitual, contemporaneidade, crítico, debate, demanda, desdobramento, dialógico, disciplina, econômico, exigir, figurar, final, horizonte, ideal, inspirar, instigar, moderno, método, organizador, pertinente, problema, razão, relevante, segar, sequência, século, tanto, tarefa, tecnologia, teoria, teórico
\end{textsample}
